% page 401 du fac-simile (image p-400 du scan)
% bloc 160.0 x 233.3 mm ; marge gauche 42.0 mm
\pgc{20.320mm}{0.000mm}{
\l{\cel{48}393.}
\l{}
\l{\cel{3}- Tubo kornatra qua quaze branchifas en la alo di ula insek-}
\l{\cel{3}ti. - II. (\sou{Bind-arto}).Salianto quan formacas, sur la dorso}
\l{\cel{3}di libro, la kordeti ye qui sutesis la kayeri. - III.(\sou{arki}-}
\l{\cel{3}\sou{tekturo}). Muluro salianta di la aristi di vulto - di la}
\l{\cel{3}kaneli di kolono, e c. - IS.}
\l{}
\l{\sou{nesto}. Mikra korbo ronda,quan la uceli konstruktas ek stipeti}
\l{\cel{3}de palio, junko, ligno, ed ek moso, lanugo, e c. por pondar}
\l{\cel{3}en olu sua ovi, kovar li, ekirigar la yuneti ed edukar ici.}
\l{\cel{3}DE.}
\l{}
\l{\sou{neta}. I. Quan nula makulo igas sordida. - II. (\sou{pri preco}).}
\l{\cel{3}De qua onu deduktis omno deduktenda : rabato, gelto, dis-}
\l{\cel{3}konto, e c.\cel{2}- EFIS.}
\l{}
\l{\sou{neurono}.\cel{2}(\sou{histol.})\cel{2}Celulo nervala, ek celulo-korpo nukleoza}
\l{\cel{3}e prolonguri, arboratraji protoplasmala ed axono. -}
\l{}
\l{\sou{neutra}.\cel{2}I.Qua ne adheras a la partiso di ulu en debato, en}
\l{\cel{3}milito. - II. (\sou{kemio}).Dicesas pri korpo ek la kombinuro}
\l{\cel{3}kemiala di du korpi (acido e bazo). - DEFIRS.}
\l{}
\l{\sou{nevo}. La filio di onua frato.}
\l{}
\l{\sou{neveo}. Strato de nivo-floki apud glaciero. EFIS.}
\l{}
\l{\sou{nevralgio}. (\sou{patol.})\cel{2}Doloro nervala. - DEFIRS.}
\l{}
\l{\sou{nevrastenio}. (\sou{patol}.)Nevroso di qua la karaktero esencala}
\l{\cel{3}esas la hipotonio neuro-muskulala e fatigeso. - DEF.}
\l{}
\l{\sou{nevrito}. (\sou{patol.}) Lezuro inflamurala o degenero di la nervi.}
\l{\cel{3}- DEF.}
\l{}
\l{\sou{nevrologio}.\cel{2}(\sou{biol.}) La cienco qua traktas la formaco, la}
\l{\cel{3}strukturo, la roli e la morbi di la sistemo nervala e,}
\l{\cel{3}specale, di la nervi. - DEFIS.}
\l{}
\l{\sou{nevroptero}. (\sou{histo}rio\cel{1}\sou{naturala)}. Insekto di qua la ali havas}
\l{\cel{3}nervaturo dispozita ret-atre. (kom ex.: libelulo).- DEFIS .}
\l{}
\l{\sou{nevroso}. (\sou{patol.})\cel{2}Stando maladatra, karakterizata da trubleso}
\l{\cel{3}nervala. - DEFIRS.}
\l{}
\l{\sou{nevuso}. (\sou{patol}.) Misformacuro di pelo, multa-kaze naskala,de}
\l{\cel{3}origino embrionala, qua havas la formo di makulo o tumoro,}
\l{\cel{3}ma qua povas eventar pos la nasko e mem che oldi. - DEF.}
\l{}
\l{\sou{+nexta}. Qua arivos kom unesma (en serio od en spaco). ED.}
\l{}
\l{ni. En semblo de la personi quin la parolanto aludas kom}
\l{\cel{3}esanta en la sama grupo kam lu.}
}
